\begin{table}[tbp]\centering\small
\caption{What the Granger Test on the Level Picks Up}\label{tab:granger}
\resizebox{\ifdim\width>\textwidth\textwidth\else\width\fi}{!}{%
\begin{tabular}{lrrr}
\toprule
 & All 12 lags & Current level (sum of lags) & 11 changes \\
\midrule
Market index, nominal & 1.9e$-$05 & 0.27 & 4.6e$-$05 \\
Market index, relative to $r^*$ & 6.0e$-$09 & 0.17 & 1.5e$-$07 \\
GS, relative to $r^*$ & 3.0e$-$11 & 0.15 & 3.7e$-$09 \\
FRB/US index, relative to $r^*$ & 5.3e$-$08 & 0.35 & 5.5e$-$07 \\
FCI-G, relative to $r^*$ & 6.3e$-$11 & 0.00 & 0.121 \\
NFCI & 2.1e$-$13 & 0.00 & 1.6e$-$04 \\
ANFCI & 1.4e$-$14 & 0.00 & 6.6e$-$04 \\
\bottomrule
\end{tabular}}
\begin{flushleft}\footnotesize\textit{Notes:} $p$-values of tests in the unemployment equation of the two-variable VAR (12 lags, 1973--2026): the 12 lags of the index jointly (the Granger test); their sum, which is the coefficient on the current level once the lags are rewritten as the level and 11 one-month changes; and the 11 changes jointly.\end{flushleft}
\end{table}
